%  Origen: 03_esquema_solucion_alcance/entrega_2/tablas/esquema_solucion_alcance.xlsx
%  Hoja "3.9 Registro de reglas de negocio".

\begin{tablalafrox}{Registro de reglas de negocio}%
  {tab:reglas-negocio}%
  {C{1.5cm} L{3.7cm} Y}%
  {\cab{ID} & \cab{Regla} & \cab{Definición operativa}}
  RNG-01 & Asignación y reserva de stock
    & La reserva se produce al confirmar el pedido, no al tomarlo ni al
      prepararlo, y se resuelve por orden de llegada. El segundo preventista
      que confirma sobre el mismo stock queda con quiebre parcial y no bloquea
      al primero. \\
  RNG-02 & Stock disponible, comprometido y en tránsito
    & Disponible es el físico menos el comprometido. El tránsito no suma hasta
      que la recepción se confirma. Sin conexión el dato es indicativo y se
      muestra con su antigüedad; los conflictos se resuelven por orden
      cronológico. \\
  RNG-03 & Política de crédito en preventa
    & Superar el límite dispara alerta; por sobre el 110~\% se bloquea el
      envío. Las excepciones las autoriza un supervisor y quedan registradas;
      sin conexión se aprueban al sincronizar. \\
  RNG-04 & Excursión térmica y disposición del producto
    & Los parámetros de tiempo y grados los define la jefa de calidad y son
      configurables. La alerta es inmediata, pero el bloqueo no es automático:
      el conductor decide la disposición según el producto. \\
  RNG-05 & Reintento de entrega y local cerrado
    & El conductor reporta local cerrado, el sistema registra el fallido y
      reagenda a la siguiente ventana del cliente. El conductor no decide. Si
      la indisponibilidad persiste, los productos vuelven al almacén. \\
  RNG-06 & Devoluciones y efecto tributario
    & La devolución se registra en la prueba de entrega, sin conexión y con
      motivo tomado de catálogo. El destino se decide en el andén de retorno y
      pasa al ERP. El documento ya emitido se resuelve con nota de crédito,
      con retención legal de seis años. \\
  RNG-07 & Control de envases retornables
    & Se lleva saldo por cliente en cuenta corriente, no por unidad
      identificada. La entrega descuenta y la devolución suma; el preventista
      recibe alerta al superar el umbral y se emite reporte diario de pérdida. \\
  RNG-08 & Cambio de precio entre toma y despacho
    & Se cobra el precio acordado por el preventista al tomar el pedido.
      Los cambios posteriores de la lista de precios no modifican el importe
      pactado. Toda discrepancia se alerta al responsable comercial y se
      corrige antes de la facturación. \\
  RNG-09 & Entrega cumplida, OTIF y fill rate
    & Una entrega parcial no cuenta como cumplida: \emph{in full} exige el
      100~\% de ítems y unidades. \emph{On time} se mide contra fecha y
      ventana. El fill rate compara unidades entregadas contra pedidas. Toda
      parcial queda como desvío auditado con causal. \\
  RNG-10 & FEFO y vida útil remanente
    & El picking sigue estrictamente el vencimiento de lote, con bitácora de
      excepción. Se alerta bajo un umbral configurable, del orden de 30 días. \\
  RNG-11 & Maestro de productos y códigos externos
    & Prevalece el maestro interno de Puelche como única fuente autorizada, con
      historial de códigos. Los códigos externos se resuelven por tabla de
      equivalencias por cadena; lo no resuelto pasa a bandeja de excepciones. \\
  RNG-12 & Sincronización y resolución de conflictos
    & Al reconectar, la sincronización es automática con reconciliación
      determinista y bitácora auditable. Durante las 14 horas de terreno el
      dato sin conexión es indicativo y el stock se valida contra el servidor
      al reconectar. \\
  RNG-13 & Secuenciación térmica y validación por zona
    & La secuencia de picking es seco, refrigerado y congelado. La
      compatibilidad térmica se valida al ubicar: un congelado no puede
      quedar fuera de zona congelada. \\
  RNG-14 & Capacidad del vehículo y carga dirigida
    & No se asigna carga que exceda el peso o el volumen nominal, ni cadena de
      frío a un vehículo sin equipo. La carga se ordena en secuencia inversa a
      la ruta y se verifica el escaneo de todos los bultos. \\
\end{tablalafrox}
